\documentclass[10pt,a4paper]{amsart}
\usepackage[margin=19mm]{geometry}
\usepackage{amsmath,amssymb,mathtools}
\usepackage[scr=rsfs,cal=euler]{mathalfa}
\usepackage{xfrac}
\usepackage[unicode,hidelinks,bookmarks=false]{hyperref}
\newcommand{\ie}{i.e.\ }
\newcommand{\cf}{cf.\ }
\newcommand{\resp}{respectively\ }
\newcommand{\Subsection}{\subsection}
\providecommand{\nobreakdash}{\nobreak\hbox{-}}
\setlength{\emergencystretch}{2em}
\multlinegap0pt
\usepackage{mathtools}
\usepackage{amssymb}
\newtheorem{proposition}{Proposition}
\renewcommand{\le}{\leqslant}
\title{Semiprojective Banach lattices}
\author{Tomasz Kania, Mariusz Niwi\'nski}
\date{Source: arXiv:2604.10624v1 (CC BY 4.0)}
\begin{document}
\maketitle

\noindent\emph{Source and licence.} Semiprojective Banach lattices. Version arXiv:2604.10624v1 (CC BY 4.0), distributed by arXiv under the Creative Commons Attribution 4.0 International licence, http://creativecommons.org/licenses/by/4.0/, as recorded in the arXiv licence metadata for this exact version and shown on the abstract page https://arxiv.org/abs/2604.10624v1. Full licence notice: \texttt{licenses/arxiv-2604.10624-CC-BY-4.0.txt}.\par\smallskip

\noindent\textit{Editorial scope.} Counted unit: the article's proposition prop:sep-quotient-lift (source0.tex lines 836--842). The article's complete original proof of it is delivered with it (source0.tex lines 844--940, one \texttt{proof} environment of 97 lines). No mathematics is written, repaired or re-derived: the statement and the proof are the article's own lines, in the article's own notation, and no retained line is substituted or re-flowed.  No further source-local context is needed: every cross-reference of every retained line resolves inside the two delivered spans.  Nothing was omitted from any retained span, so the package carries no omission record.  No retained line contains a citation, so the wrapper carries no bibliography and no bibliography entry.  No retained line contains a figure environment, an image-inclusion command or drawing code, so no image is delivered and the package declares no asset dependency.  Editorial formatting only, in addition: an AMS \texttt{amsart} preamble that restates the article's own declarations and macros used by the retained lines, namely the environments \texttt{proposition} on separate counters, together with the standard packages those lines need.  The retained environments print in this document as Proposition 1 in order of appearance, whereas the article prints the same items with the numbering of its own full document.  The complete native source of the article as distributed in the arXiv e-print for this exact version is bundled unchanged as \texttt{sources/198283/source0.tex}.
\begin{proposition}[Lifting through quotients of separable Banach lattices]\label{prop:sep-quotient-lift}
Let $A$ be a semiprojective Banach lattice, let $X$ be a separable Banach lattice,
let $J\subseteq X$ be a closed ideal, and let
$T\colon A\to X/J$ be a contractive lattice homomorphism.
Then there exists a lattice homomorphism $\widehat T\colon A\to X$ such that
$Q\circ \widehat T=T$, where $Q\colon X\to X/J$ is the quotient map.
\end{proposition}

\begin{proof}
Consider
\[
B:=\Bigl\{f\in C([0,1],X): f(0)=0
\text{ and } Q(f(t))=t\,Q(f(1)) \text{ for all } t\in[0,1]\Bigr\},
\]
equipped with the supremum norm and the pointwise order.
Then \(B\) is a closed sublattice of \(C([0,1],X)\): closedness follows because the
defining conditions are preserved under uniform limits, and if \(f\in B\), then
\[
Q(|f(t)|)=|Q(f(t))|=t\,|Q(f(1))|=t\,Q(|f(1)|)
\qquad (t\in[0,1]),
\]
so \(|f|\in B\).

Since $X$ is separable and $[0,1]$ is compact metric, $C([0,1],X)$ is separable,
hence so is $B$.

Define
\[
\Pi\colon B\to X/J,\qquad \Pi(f):=Q(f(1)).
\]
Then $\Pi$ is a contractive lattice homomorphism.
It is surjective: given $y\in X/J$, choose $x\in X$ with $Qx=y$ and set
$f_x(t):=tx$; then $f_x\in B$ and $\Pi(f_x)=y$.

Its kernel is
\[
I=\{f\in C([0,1],J): f(0)=0\}.
\]
For each $n\in\mathbb N$, let
\[
I_n:=\{f\in I: f|_{[0,1/n]}=0\}.
\]
Then each \(I_n\) is a closed ideal in \(B\), the sequence \((I_n)\) is increasing,
and \(I=\overline{\bigcup_n I_n}\). Indeed, closedness and monotonicity are clear. If
\(g\in I_n\) and \(f\in B\) satisfy \(|f|\le |g|\), then \(f(t)=0\) for \(t\in[0,1/n]\),
and since \(g(t)\in J\) and \(J\) is an ideal in \(X\), also \(f(t)\in J\) for every
\(t\in[0,1]\); hence \(f\in I_n\).
Indeed, if $f\in I$, choose $u_n\in C([0,1])$ with
$0\le u_n\le 1$, $u_n|_{[0,1/n]}=0$, and $u_n|_{[2/n,1]}=1$.
Then $u_nf\in I_n$ and
\[
\|f-u_nf\|_\infty\le \sup_{0\le t\le 2/n}\|f(t)\|\xrightarrow[n\to\infty]{}0,
\]
because $f(0)=0$.

The map $\Pi$ induces a lattice isomorphism
\[
\widetilde\Pi\colon B/I\to X/J.
\]
Moreover, $\widetilde\Pi$ is isometric: contractivity is clear, while for
$y\in X/J$ and $\eta>0$ one may choose $x\in X$ with $Qx=y$ and
$\|x\|\le \|y\|+\eta$; then $f_x(t)=tx$ belongs to $B$ and
\[
\|\widetilde\Pi^{-1}(y)\|\le \|f_x\|_\infty=\|x\|\le \|y\|+\eta.
\]
Letting $\eta\downarrow 0$ yields $\|\widetilde\Pi^{-1}(y)\|\le \|y\|$.

Now set
\[
\phi:=\widetilde\Pi^{-1}\circ T\colon A\to B/I.
\]
Since $A$ is semiprojective and $B$ is separable, there exist $n\in\mathbb N$
and, for instance with $\varepsilon=1$, a lattice homomorphism
\[
\psi\colon A\to B/I_n
\]
such that
\[
\pi_n\circ \psi=\phi,
\]
where $\pi_n\colon B/I_n\to B/I$ is the canonical quotient map.

Define
\[
E_n\colon B/I_n\to X,\qquad E_n([f]):=n\,f(1/n).
\]
This is well-defined because every element of $I_n$ vanishes at $1/n$.
It is a lattice homomorphism, since point evaluation and multiplication by a
positive scalar preserve lattice operations.
For $[f]\in B/I_n$ we have
\[
Q(E_n([f]))
=n\,Q(f(1/n))
=n\cdot \frac1n\,Q(f(1))
=\widetilde\Pi(\pi_n([f])).
\]
Hence
\[
Q\circ E_n\circ \psi
=\widetilde\Pi\circ \pi_n\circ \psi
=\widetilde\Pi\circ \phi
=T.
\]
Therefore $\widehat T:=E_n\circ \psi$ is the desired lift.
\end{proof}

\end{document}
